% !PRINTED 272
\weberpage{272}
%<<CONT
f\"ur die man aus (2) die zusammengeh\"origen Vertauschungen erh\"alt:
\weq{12}{\begin{array}{lll}
\omega, & A, & B,\\
-\dfrac{1}{\omega}, & \varrho^2 A, & \varrho B,\\
\omega+1, & e^{\frac{(n+1)\pi i}{6}}\sigma^2 A, & e^{-\frac{(n+1)\pi i}{6}}\sigma B,
\end{array}}
und hieraus leitet man in der gleichen Weise die Gleichungen ab:
\weq{13}{\begin{aligned}
n&=3, & A&=0,\\
n&=11, & A&=1,\\
n&=19, & A^5-7A^2-B&=0.
\end{aligned}}
3. Ist $n \equiv 1$ (mod 4), so kann man ebenso verfahren mit den Funktionen
\weq{14}{\begin{aligned}
8A &= u^4v^4 - u_1^4v_1^4 - u_2^4v_2^4,\\
B &= u^4v^4u_1^4v_1^4 + u^4v^4u_2^4v_2^4 - u_1^4v_1^4u_2^4v_2^4,
\end{aligned}}
f\"ur die man die zusammengeh\"origen Vertauschungen erh\"alt:
\weq{15}{\begin{array}{lll}
\omega, & A, & B,\\
-\dfrac{1}{\omega}, & \varrho\,A, & \varrho^2 B,\\
\omega+1, & e^{\frac{(n+1)\pi i}{3}}\sigma\,A, & e^{-\frac{(n+1)\pi i}{3}}\sigma^2 B.
\end{array}}
Nur der erste Fall $n=5$ f\"uhrt hier zu einem einfachen Resultat:
\weq{16}{n=5, \quad A=1.}
F\"ur $n=5$ l\"a\ss{}t sich noch eine einfachere Form der Transformationsgleichung gewinnen.

Wenn wir n\"amlich auf die drei Funktionen
\weq{17}{w = \frac{u_1^2v_2^2+u_2^2v_1^2}{uv}, \quad w_1 = \frac{u_2^2v^2-v_2^2u^2}{u_1v_1}, \quad w_2 = \frac{u_1^2v^2-v_1^2u}{u_2v_2}} % [sic] Source prints u without the square in the w_2 numerator.
die Substitutionen $\bigl(\omega, -\frac{1}{\omega}\bigr)$, $(\omega, \omega+1)$ anwenden, so ergeben sich unter der Voraussetzung $n \equiv 5$ (mod 8) die zusammengeh\"origen Vertauschungen:
\[
\begin{array}{llll}
\omega, & w, & w_1, & w_2,\\
-\dfrac{1}{\omega}, & w, & w_2, & w_1,\\
\omega+1, & e^{-\frac{\pi i}{24}(n-5)}w_1, & e^{-\frac{\pi i}{24}(n-5)}w, & e^{\frac{\pi i}{12}(n-5)}w_2.
\end{array}
\]
